\section{Limitations and Future Work}
\label{sec:limitations}

\paragraph{Limitations.}
(1) \textbf{Domain.} Project-Utopia is a colony simulation in a
2-D tile world. Generalisation of stale coherence as a phenomenon
to other long-horizon agent domains --- city-traffic control,
hospital scheduling, supply-chain orchestration, dialogue-heavy
assistants, robotics primitives --- is left as future work. The
decay profiles we measure may differ in dialogue-heavy or
robotics-heavy domains; we characterise one regime in depth rather
than many regimes shallowly.
(2) \textbf{Sample size.} The main matrix is 5 seeds $\times$ 3
models $\times$ 3 scenarios $=$ 15 runs per
(model, scenario) cell for E1/E2, and 5 seeds per
context-budget cell for E3. We adopt
\citet{Bowyer2025DontUseClt}'s Beta-Binomial methodology
appropriate for small-$N$, but acknowledge that doubling seed
count would tighten credible intervals and may surface effects
below our current detection threshold.
(3) \textbf{LLM API stationarity.} Anthropic, OpenAI, and the
open-weight hosts we use all rev models and adjust safety filters
on inference; our model-snapshot pinning slows but does not
eliminate this. The RecordReplayCache lets reviewers replay our
exact prompts-and-responses transcript bit-identically, but
re-running with the latest production endpoint is not guaranteed
to reproduce.
(4) \textbf{Anchor count and design.} Stale coherence may scale
non-trivially with anchor count or anchor complexity. We use a
fixed anchor count of 5 anchors per episode, each carrying a
paired (verbal token, implicit-goal) structure. Other designs ---
single anchor, 10 anchors, contradictory anchors, dynamically
introduced anchors --- are out of scope for this paper and are
the most immediate ablation we plan post-submission.
(5) \textbf{Partial observability.} Project-Utopia currently
exposes the full world summary to all four channels. A more
realistic deployment would give each channel only its proper
observation slice. We expect partial observability to
\emph{increase} $\Delta$ (less verbal echo per probe, similar
action grounding) and treat this as a useful future ablation,
not as a threat to the present results.
(6) \textbf{Inference-time only.} We measure inference-time
stale coherence. Whether RLHF / DPO with $\Delta$ or the dual
probe as part of the reward signal reduces stale coherence is
an open question that requires training-time access we do not
have for the frontier closed-weight models in our menu.
(7) \textbf{Anchor protocol design choice.} Our anchors carry
explicit verbal tokens that we then probe for. An alternative
protocol, in which anchors carry only implicit context to be
inferred, might surface different decay profiles. The current
choice maximises measurability (the verbal probe has a
well-defined target string); the alternative would maximise
ecological validity (production anchors are rarely re-stated
verbatim).
(8) \textbf{Empirical evidence pending.} At submission, no
real-LLM runs are populated for E1, E2, or E3. The pipeline is
verified end-to-end on fallback NDJSON: schema validity, anchor
injection, dual-probe sampling, statistical machinery, and
figure rendering all run through the released CLI. Cell-level
real-LLM signal is reserved for the camera-ready once API
access is provisioned.

\paragraph{Future work.}
Camera-ready will populate E1, E2, and E3 with the 3-family model
menu in \S\ref{sec:exp:setup}
(Claude-Sonnet-4.6, GPT-5-mini, Llama-3.1-8B-Instruct).
Subsequent work extends the dual-probe methodology to other
long-horizon agent domains (dialogue assistants, robotics
primitives) and explores whether RLHF or DPO on $\Delta$
reduces stale coherence at training time.

\paragraph{What we explicitly do \emph{not} promise.}
Project-Utopia does not claim to test ``general intelligence'',
linguistic competence, or commonsense reasoning. It tests
\emph{stale coherence under multi-cadence structured output} ---
specifically, whether an LLM continues to act on an anchor it
continues to verbally mention, over an 8-sim-hour episode. Models
that do well on Project-Utopia are not necessarily strong on
GPQA / MMLU / HumanEval; conversely, GPQA-strong models can fail
stale-coherence tests for reasons orthogonal to ``reasoning''
(e.g., the verbal-echo channel and the action-grounding channel
of the model are decoupled in ways the reasoning benchmarks do
not surface).
